% !TEX encoding = UTF-8
% Общий преамбул для статей и глав диссертации (LuaLaTeX)

% Поддержка русского языка
\usepackage[english,russian]{babel}
\usepackage[autostyle=true,english=american,russian]{csquotes}

% Шрифты (Times New Roman по ГОСТ)
\usepackage{fontspec}
\setmainfont{Times New Roman}
\defaultfontfeatures{Ligatures=TeX}

% Класс документа НЕ указывается здесь — только в main.tex

% Поля по ГОСТ Р 7.0.11–2011
\usepackage{geometry}
\geometry{
  a4paper,
  top=20mm,
  bottom=20mm,
  left=30mm,
  right=15mm,
  nohead
}

% Рисунки, подписи, плавающие объекты
\usepackage{graphicx}
\usepackage{caption}
\usepackage{float}
\captionsetup{
  labelsep=endash,
  justification=justified,
  singlelinecheck=false
}

% Списки
\usepackage{enumitem}
\setlist[itemize]{left=0pt}
\setlist[enumerate]{left=0pt}

% Математика
\usepackage{amsmath, amssymb}

% Библиография (BibLaTeX + Biber)
\usepackage[
  backend=biber,
  style=gost-numeric,
  sorting=none,
  doi=false,
  isbn=false,
  url=false
]{biblatex}
\addbibresource{../references.bib} % относительно articles/xxx/

% Гиперссылки без цвета (по ГОСТ)
\usepackage[hidelinks]{hyperref}

% Форматирование заголовков (без точки в конце)
\usepackage{titlesec}
\titleformat{\section}{\normalfont\large\bfseries}{\thesection}{1em}{}
\titleformat{\subsection}{\normalfont\bfseries}{\thesubsection}{1em}{}

% Укажем, где искать рисунки — относительно main.tex
\graphicspath{{./}{images/}}

% Типографика
\tolerance=1000
\emergencystretch=10pt